% !TeX encoding = UTF-8
% !TeX program = xelatex
% !TeX spellcheck = en_US

\documentclass[degree=master]{thuthesis}
  % 学位 degree:
  %   doctor | master | bachelor | postdoc
  % 学位类型 degree-type:
  %   academic（默认）| professional
  % 语言 language
  %   chinese（默认）| english
  % 字体库 fontset
  %   windows | mac | fandol | ubuntu
  % 建议终版使用 Windows 平台的字体编译


% 论文基本配置，加载宏包等全局配置
% !TeX root = ./thuthesis-example.tex

% 论文基本信息配置

\thusetup{
  %******************************
  % 注意：
  %   1. 配置里面不要出现空行
  %   2. 不需要的配置信息可以删除
  %   3. 建议先阅读文档中所有关于选项的说明
  %******************************
  %
  % 输出格式
  %   选择打印版（print）或用于提交的电子版（electronic），前者会插入空白页以便直接双面打印
  %
  output = print,
  %
  % 标题
  %   可使用“\\”命令手动控制换行
  %
  title  = {清华大学学位论文 \LaTeX{} 模板\\使用示例文档 v\version},
  title* = {An Introduction to \LaTeX{} Thesis Template of Tsinghua
            University v\version},
  %
  % 学科门类
  %   1. 学术型
  %      - 中文
  %        需注明所属的学科门类，例如：
  %        哲学、经济学、法学、教育学、文学、历史学、理学、工学、农学、医学、
  %        军事学、管理学、艺术学
  %      - 英文
  %        博士：Doctor of Philosophy
  %        硕士：
  %          哲学、文学、历史学、法学、教育学、艺术学门类，公共管理学科
  %          填写“Master of Arts“，其它填写“Master of Science”
  %   2. 专业型
  %      直接填写专业学位的名称，例如：
  %      教育博士、工程硕士等
  %      Doctor of Education, Master of Engineering
  %   3. 本科生不需要填写
  %
  degree-category  = {工学硕士},
  degree-category* = {Master of Science},
  %
  % 培养单位
  %   填写所属院系的全名
  %
  department = {计算机科学与技术系},
  %
  % 学科
  %   1. 研究生学术型学位，获得一级学科授权的学科填写一级学科名称，其他填写二级学科名称
  %   2. 本科生填写专业名称，第二学位论文需标注“（第二学位）”
  %
  discipline  = {计算机科学与技术},
  discipline* = {Computer Science and Technology},
  %
  % 专业领域
  %   1. 设置专业领域的专业学位类别，填写相应专业领域名称
  %   2. 2019 级及之前工程硕士学位论文，在 `engineering-field` 填写相应工程领域名称
  %   3. 其他专业学位类别的学位论文无需此信息
  %
  % professional-field  = {计算机技术},
  % professional-field* = {Computer Technology},
  %
  % 姓名
  %
  author  = {薛瑞尼},
  author* = {Xue Ruini},
  %
  % 指导教师
  %   中文姓名和职称之间以英文逗号“,”分开，下同
  %
  supervisor  = {郑纬民, 教授},
  supervisor* = {Professor Zheng Weimin},
  %
  % 副指导教师
  %
  associate-supervisor  = {陈文光, 教授},
  associate-supervisor* = {Professor Chen Wenguang},
  %
  % 联合指导教师
  %
  % co-supervisor  = {某某某, 教授},
  % co-supervisor* = {Professor Mou Moumou},
  %
  % 日期
  %   使用 ISO 格式；默认为当前时间
  %
  % date = {2019-07-07},
  %
  % 是否在中文封面后的空白页生成书脊（默认 false）
  %
  include-spine = false,
  %
  % 密级和年限
  %   秘密, 机密, 绝密
  %
  % secret-level = {秘密},
  % secret-year  = {10},
  %
  % 博士后专有部分
  %
  % clc                = {分类号},
  % udc                = {UDC},
  % id                 = {编号},
  % discipline-level-1 = {计算机科学与技术},  % 流动站（一级学科）名称
  % discipline-level-2 = {系统结构},          % 专业（二级学科）名称
  % start-date         = {2011-07-01},        % 研究工作起始时间
}

% 载入所需的宏包

% 定理类环境宏包
\usepackage{amsthm}
% 也可以使用 ntheorem
% \usepackage[amsmath,thmmarks,hyperref]{ntheorem}

\thusetup{
  %
  % 数学字体
  % math-style = GB,  % GB | ISO | TeX
  math-font  = xits,  % stix | xits | libertinus
}

% 可以使用 nomencl 生成符号和缩略语说明
% \usepackage{nomencl}
% \makenomenclature

% 表格加脚注
\usepackage{threeparttable}

% 表格中支持跨行
\usepackage{multirow}

% 固定宽度的表格。
% \usepackage{tabularx}

% 跨页表格
\usepackage{longtable}

% 算法
\usepackage{algorithm}
\usepackage{algorithmic}

% 量和单位
\usepackage{siunitx}

% 参考文献使用 BibTeX + natbib 宏包
% 顺序编码制
\usepackage[sort]{natbib}
\bibliographystyle{thuthesis-numeric}

% 著者-出版年制
% \usepackage{natbib}
% \bibliographystyle{thuthesis-author-year}

% 本科生参考文献的著录格式
% \usepackage[sort]{natbib}
% \bibliographystyle{thuthesis-bachelor}

% 参考文献使用 BibLaTeX 宏包
% \usepackage[style=thuthesis-numeric]{biblatex}
% \usepackage[style=thuthesis-author-year]{biblatex}
% \usepackage[style=apa]{biblatex}
% \usepackage[style=mla-new]{biblatex}
% 声明 BibLaTeX 的数据库
% \addbibresource{ref/refs.bib}

% 定义所有的图片文件在 figures 子目录下
\graphicspath{{figures/}}

% 数学命令
\makeatletter
\newcommand\dif{%  % 微分符号
  \mathop{}\!%
  \ifthu@math@style@TeX
    d%
  \else
    \mathrm{d}%
  \fi
}
\makeatother

\usepackage{listings}

\lstset{
	columns=fixed,
  basicstyle=\small,
	numbers=left,                                        % 在左侧显示行号
	numberstyle=\tiny\color{gray},                       % 设定行号格式
	frame=none,                                          % 不显示背景边框
	backgroundcolor=\color[RGB]{245,245,244},            % 设定背景颜色
	keywordstyle=\color[RGB]{40,40,255},                 % 设定关键字颜色
	numberstyle=\footnotesize\color{darkgray},
	commentstyle=\it\color[RGB]{0,96,96},                % 设置代码注释的格式
	stringstyle=\rmfamily\slshape\color[RGB]{128,0,0},   % 设置字符串格式
	showstringspaces=false,                              % 不显示字符串中的空格
	breaklines=true,
	language=java,                                        % 设置语言
}

\renewcommand{\lstlistingname}{算法}% Listing -> Algorithm
% hyperref 宏包在最后调用
\usepackage{hyperref}



\begin{document}

% 封面
% \maketitle

% 学位论文指导小组、公开评阅人和答辩委员会名单
% 本科生不需要
% % !TeX root = ../thuthesis-example.tex

\begin{committee}[name={学位论文指导小组、公开评阅人和答辩委员会名单}]

  \newcolumntype{C}[1]{@{}>{\centering\arraybackslash}p{#1}}

  \section*{指导小组名单}

  \begin{center}
    \begin{tabular}{C{3cm}C{3cm}C{9cm}@{}}
      李XX & 教授     & 清华大学 \\
      王XX & 副教授   & 清华大学 \\
      张XX & 助理教授 & 清华大学 \\
    \end{tabular}
  \end{center}


  \section*{公开评阅人名单}

  \begin{center}
    \begin{tabular}{C{3cm}C{3cm}C{9cm}@{}}
      刘XX & 教授   & 清华大学                    \\
      陈XX & 副教授 & XXXX大学                    \\
      杨XX & 研究员 & 中国XXXX科学院XXXXXXX研究所 \\
    \end{tabular}
  \end{center}


  \section*{答辩委员会名单}

  \begin{center}
    \begin{tabular}{C{2.75cm}C{2.98cm}C{4.63cm}C{4.63cm}@{}}
      主席 & 赵XX                  & 教授                    & 清华大学       \\
      委员 & 刘XX                  & 教授                    & 清华大学       \\
          & \multirow{2}{*}{杨XX} & \multirow{2}{*}{研究员} & 中国XXXX科学院 \\
          &                       &                         & XXXXXXX研究所  \\
          & 黄XX                  & 教授                    & XXXX大学       \\
          & 周XX                  & 副教授                  & XXXX大学       \\
      秘书 & 吴XX                  & 助理研究员              & 清华大学       \\
    \end{tabular}
  \end{center}

\end{committee}



% 也可以导入 Word 版转的 PDF 文件
% \begin{committee}[file=figures/committee.pdf]
% \end{committee}


% 使用授权的说明
% \copyrightpage
% 将签字扫描后授权文件 scan-copyright.pdf 替换原始页面
% \copyrightpage[file=scan-copyright.pdf]

\frontmatter
% !TeX root = ../thuthesis-example.tex

% 中英文摘要和关键字

\begin{abstract}
	本协议基于AES-GCM和x25519实现了类似TSL1.3协议的功能，在仅两方知道密码又不要求证书的情况下，仅需 1-RTT 做到内容安全、端点可靠认证，无法通过主动探测、重放攻击、机器学习等方法区分本协议与真实TLS1.3。并且为了提高机器和网络性能，在不影响安全的前提下，可以选择性地加密和使用网络多路复用。

	% 关键词用“英文逗号”分隔，输出时会自动处理为正确的分隔符
	\thusetup{
		keywords = {TLS1.3, 信道安全, 伪造TLS},
	}
\end{abstract}

\begin{abstract*}
	This protocol implements functions similar to TSL1.3 protocol based on AES-GCM and x25519. When only two parties know the password and do not require a certificate, only 1-RTT is required to achieve content security and reliable endpoint authentication. It is impossible to distinguish this protocol from real TLS1.3 through active detection, Replay attack, machine learning and other methods. And in order to improve machine and network performance, selective encryption and network multiplexing can be used without affecting security.

	% Use comma as separator when inputting
	\thusetup{
		keywords* = {TLS1.3, Channel Security},
	}
\end{abstract*}


% 目录
\tableofcontents

% 插图和附表清单
% 本科生的插图索引和表格索引需要移至正文之后、参考文献前
% \listoffiguresandtables  % 插图和附表清单（仅限研究生）
\listoffigures           % 插图清单
\listoftables            % 附表清单

% 符号对照表
% !TeX root = ../thuthesis-example.tex

\begin{denotation}[3cm]
  \item[Server] 服务端
  \item[Client] 客户端
  \item[TLS] 传输层安全性协议(Transport Layer Security)
  \item[ECDHE] 椭圆曲线迪菲-赫尔曼金钥交换 (Elliptic-curve Diffie–Hellman)
\end{denotation}



% 也可以使用 nomencl 宏包，需要在导言区
% \usepackage{nomencl}
% \makenomenclature

% 在这里输出符号说明
% \printnomenclature[3cm]

% 在正文中的任意为都可以标题
% \nomenclature{PI}{聚酰亚胺}
% \nomenclature{MPI}{聚酰亚胺模型化合物，N-苯基邻苯酰亚胺}
% \nomenclature{PBI}{聚苯并咪唑}
% \nomenclature{MPBI}{聚苯并咪唑模型化合物，N-苯基苯并咪唑}
% \nomenclature{PY}{聚吡咙}
% \nomenclature{PMDA-BDA}{均苯四酸二酐与联苯四胺合成的聚吡咙薄膜}
% \nomenclature{MPY}{聚吡咙模型化合物}
% \nomenclature{As-PPT}{聚苯基不对称三嗪}
% \nomenclature{MAsPPT}{聚苯基不对称三嗪单模型化合物，3,5,6-三苯基-1,2,4-三嗪}
% \nomenclature{DMAsPPT}{聚苯基不对称三嗪双模型化合物（水解实验模型化合物）}
% \nomenclature{S-PPT}{聚苯基对称三嗪}
% \nomenclature{MSPPT}{聚苯基对称三嗪模型化合物，2,4,6-三苯基-1,3,5-三嗪}
% \nomenclature{PPQ}{聚苯基喹噁啉}
% \nomenclature{MPPQ}{聚苯基喹噁啉模型化合物，3,4-二苯基苯并二嗪}
% \nomenclature{HMPI}{聚酰亚胺模型化合物的质子化产物}
% \nomenclature{HMPY}{聚吡咙模型化合物的质子化产物}
% \nomenclature{HMPBI}{聚苯并咪唑模型化合物的质子化产物}
% \nomenclature{HMAsPPT}{聚苯基不对称三嗪模型化合物的质子化产物}
% \nomenclature{HMSPPT}{聚苯基对称三嗪模型化合物的质子化产物}
% \nomenclature{HMPPQ}{聚苯基喹噁啉模型化合物的质子化产物}
% \nomenclature{PDT}{热分解温度}
% \nomenclature{HPLC}{高效液相色谱（High Performance Liquid Chromatography）}
% \nomenclature{HPCE}{高效毛细管电泳色谱（High Performance Capillary lectrophoresis）}
% \nomenclature{LC-MS}{液相色谱-质谱联用（Liquid chromatography-Mass Spectrum）}
% \nomenclature{TIC}{总离子浓度（Total Ion Content）}
% \nomenclature{\textit{ab initio}}{基于第一原理的量子化学计算方法，常称从头算法}
% \nomenclature{DFT}{密度泛函理论（Density Functional Theory）}
% \nomenclature{$E_a$}{化学反应的活化能（Activation Energy）}
% \nomenclature{ZPE}{零点振动能（Zero Vibration Energy）}
% \nomenclature{PES}{势能面（Potential Energy Surface）}
% \nomenclature{TS}{过渡态（Transition State）}
% \nomenclature{TST}{过渡态理论（Transition State Theory）}
% \nomenclature{$\increment G^\neq$}{活化自由能（Activation Free Energy）}
% \nomenclature{$\kappa$}{传输系数（Transmission Coefficient）}
% \nomenclature{IRC}{内禀反应坐标（Intrinsic Reaction Coordinates）}
% \nomenclature{$\nu_i$}{虚频（Imaginary Frequency）}
% \nomenclature{ONIOM}{分层算法（Our own N-layered Integrated molecular Orbital and molecular Mechanics）}
% \nomenclature{SCF}{自洽场（Self-Consistent Field）}
% \nomenclature{SCRF}{自洽反应场（Self-Consistent Reaction Field）}



% 正文部分
\mainmatter
% !TeX root = ../thuthesis-example.tex
\renewcommand{\today}{\number\year 年 \number\month 月 \number\day 日}

\chapter{协议现状}

最新版本为 V1，最后更新于 \today，处于征集讨论阶段，无协议其他变种。唯一官网：\href{https://github.com/JimmyHuang454/JLS}{https://github.com/JimmyHuang454/JLS}。文档语言：中文简体。

目前支持 JLS 协议的软件工具有：
\begin{itemize}
	\item \href{https://github.com/JimmyHuang454/RRS}{https://github.com/JimmyHuang454/RRS}
\end{itemize}



\section{Copyright}
使用最宽松、最自由的版权协议 WTFPL。

\section{贡献与反馈}
请移步到github官网，根据需要提出问题或贡献。所有贡献代码都应符合 WTFPL。

\section{Change Log}
\subsection{V2}
2023 年 6 月 30 日

更新Application Data 加解密时所使用的IV，使其更加随机。


\subsection{V1}
2023 年 6 月 24 日，初始版本

% !TeX root = ../thuthesis-example.tex

\chapter{Specification V2}
本章描述了 JLS 协议第 2 版的具体内容（不含正确性和安全性验证），另付一个基于 Dart 编程语言的具体实现。

\section{约定}
\begin{itemize}
	\item 数组相加有先后顺序，如无特殊标注，一律以大端字节序存储
	\item 如无特殊标注，本文所指TLS均为TLS版本1.3\cite{rescorla2018transport}
\end{itemize}

\section{背景}

对代理协议内容加密和认证是一项广泛的需求，基于应用层的代理协议如 HTTP、SOCK5 等都不支持内容加密\footnote{虽然有认证，但形同虚设}，因此出现了如 ShadowSocket 等加密代理协议，随着 TLS 的广泛应用，为了解决维护难、缺少安全性验证、主动探测等问题，又出现了以 TLS 为基础的协议，如 Trojan 协议。

Trojan 可以说是一个完美的协议，但必须要求购买域名，申请域名证书。在某些极端环境，仅允许特定域名通行，Trojan 失效。但预估 IP 白名单模式\footnote{仅限指定IP通行的模式}很快来临，届时我们只能祈祷。


\section{基本原理}

首先考虑一个简化模型：
现有一个客户端、一个服务端和一个仅双方知道的一个密码PWD；如何保证：\newline
\begin{itemize}
	\item 1. 客户端和服务端之间的数据原文没被篡改
	\item 2. 不能被除了客户端和服务端以外的第三方解密
	\item 3. 客户端（服务端）能够识别接收到的数据是否真的来自服务端（客户端）
	\item 4. 收到的通讯数据是否已经接收过（即防止重放攻击）\newline
\end{itemize}

为了实现第 1、2 、3 点，可以采用 AES-GCM 对称加密算法；AES 保证了第 2 点，GCM 保证了第 1 和第 3 点；

为了实现第 4 点，我们可以在 AES-GCM 的基础上往数据原文中加上一个时间戳，但是这种方式有延迟，比如说 vmess 协议，只要攻击速度够快，90 秒内可以无限让客户端或者服务端重复接收已经接收过的数据。

还有一种方法是让客户端和服务端之间进行一次握手，客户端先生成一个随机数 N1，通过 AES-GCM 加密发送 N1，服务端接收后可以解密出来，再生成一个随机数 N2，同样地加密发送给客户端；

完成握手后，就约定以后发送的数据的秘钥都是N=N1+N2+PWD，而且每发一个数据包都要在原文中附上序号（第一个包、第二包……第 x 个包）。这种会增加一次 RTT，但我们可以稍微修改一下 TLS 协议，把其中的随机数（random）换成我们给出的随机数 N1 和 N2，达到交换随机数的目的。

这种方法没有延迟，两个随机数共同组成本次对话的ID，所以该随机数应该要足够随机才能保证能防止重放攻击（Replay Attacks）。

根据 GCM 的认证功能，如果不知道 PWD，就无法解出密文，所以我们就可以检测出客户端/服务端是否假冒，从而把数据转发到我们想要的任何地址。

但为了避免服务器被探测到，只能支持全程加密的 TLS1.3；同样的，如果服务端检测到不是 TLS1.3，直接转发数据，不作处理。



\section{准备}
Client 和 Server 都必须准备一个仅双方知道的密码 PWD 和一个随机数 IV。PWD和IV长度无限制，但都不应该小于 32 字节并应该足够随机。

Server 端必须准备一个伪装站（Fallback Website）地址，伪装站应该优先选取延迟低的、使用 TLS1.3 的、使用 HTTP2 的、不能是HTTP3、DDOS 防御没那么敏感的。

强烈建议至少每1年更换一次密码和IV，最好每三个月更换一次。

\section{Handshake（握手）}
JLS的握手包跟TLS所规定的数据结构和时序完全一致。

\subsection{Client Hello 包}
功能：向服务器发送建立安全信道的请求

\begin{itemize}
	\item 如果 Client 支持 TLS1.3，那么在 Client Hello 的supported\_version拓展必须包含TLS1.3字段。JLS 是建立在 TLS1.3的基础上，所以 JLS 的 Client Hello 包必须要有此字段

	\item Client Hello 中的random字段，必须是通过算法\ref{list:clientFakeRondom}所生成的FakeRandom，通过代码\ref{list:clientFakeRondom}还会生成一个随机数 N1，N1 必须要足够随机，用于防止已知明文攻击（Known Plaintext Attack）

	\item Client Hello 中的session ID字段，应该是随机的

	\item Client Hello 包必须包含key\_share字段和交换秘钥所用算法。Client 必须自行生成正确的key\_share和所用算法（即supported\_groups）发送给 Server

	\item 开发者可以在不修改或缺少 JLS 必要信息的情况下伪造Client Hello头特征，比如说 chrome浏览器的cipher suits列表等，用于伪造 Client Hello 指纹。需要特别注意的是：如果开发者错误地伪造 Client Hello，也就是JLS能够正常使用，但 TLS 不能使用的情况下，会导致明显的特征，比如说 Client Hello 缺少了cipher suits等必要信息，TLS Server 应该发送 HelloRetryRequest，但JLS 认为这是合法 Client，所以是不会发送 HelloRetryRequest

	\item Client Hello 不应该使用early data或pre\_shared\_key，虽然可以做到 0-RTT，但是会导致重放攻击等安全问题

	\item 使用padding extension 是为了希望避免发送长度为256-511字节的Client hello\cite{langley2015rfc}，因为有不少TLS server拒绝接收长度小于512字节的Client hello，因此开发者在伪装Client hello 时，应该尽可能使用长度（不含server\_name扩展）大于512字节的Client hello
\end{itemize}

\lstinputlisting[caption=生成Client Hello,label={list:clientFakeRondom}] {data/code/build_client_hello.java}


\subsection{Server 验证}
功能：验证收到的握手是否来自有效的Client
\begin{itemize}
	\item Server 必须通过算法\ref{list:serverCheckClientHello}判断是否为有效 Client

	\item 如果不是有效 Client，直接把接收到的所有数据转发到伪装站，不作其他处理

	\item Server FakeRandom的后8字节不能是\[ 44,4F,57,4E,47,52,44,01 \]和\[ 44,4F,57,4E,47,52,44,00 \]，如果是，应该不断重新生成一个Server FakeRandom ，直到合规后再发送 Server Hello

\end{itemize}
\lstinputlisting[caption=检查是否有效Client,label={list:serverCheckClientHello}] {data/code/server_check_client_hello.java}

\subsection{发送Server Hello和证书}
功能：发送Server Hello，并发送伪装站证书
\begin{itemize}
	\item 首先要确认是否有效Client，然后通过算法\ref{list:buildServerHello}生成Server Hello.

	\item 如果是有效 Client，则可以得出来自Client的随机数N1，然后Server要生成自己的随机数N2和Server FakeRandom

	\item Server 必须要根据Client的key\_share和supported\_groups得出共同秘钥S1，并把S1也作为最终秘钥之一来加密数据，以保证前向安全性

	\item 因为 Server 已经验证了 Client 的有效性，所以 Server 证书可以传输随意内容，但包长度应该要伪装站返回的证书一致，建议开发者在软件初始化时获取伪装站的真实证书包；为了方便，返回差不多长度的 Server 证书包也是被允许的，因为这不会影响安全性，但可能成为特征；Client 无需验证该证书是否有效

	\item Session ID应与Client Hello中的一致
	\item support\_group 应只使用 x25519
\end{itemize}
\lstinputlisting[caption=生成Server Hello,label={list:buildServerHello}] {data/code/build_server_hello.java}

\subsection{Client 验证}
功能：验证是否有效Server
\begin{itemize}
	\item 根据算法\ref{list:checkServer}得出验证是否有效Server，并且得出共同秘钥S1。如果不是有效 Server，则 Client 完全按照 TLS1.3 流程处理，即验证证书，并协商出 TLS 最终随机数（秘钥），最后发送合规的 http 请求即可

	\item 如果是有效 Server，则无需验证来自 Server 的证书。
\end{itemize}
\lstinputlisting[caption=验证 server,label={list:checkServer}] {data/code/client_check_server.java}

\section{Application Data}
使用经过x25519算法得出共享密钥S1，最后得出的秘钥finalPWD = PWD + S1 + N1 + N2，通过AES\_GCM\_256加密发送数据。包结构与TLS一致。
Server和Client都必须各自维护一个自增ID，用于记录已接收和已发送包数量，按照算法\ref{list:encryptData}得出实际发送数据。目前GCM的MAC的长度为默认的16字节，每一个 Application Data 都必须在密文前加上MAC。如果Client或Server收到验证失败的Application data，必要按照正常TLS流程处理。
\lstinputlisting[caption=加密数据,label={list:encryptData}] {data/code/data_encrypt.java}

\lstinputlisting[caption=解密数据,label={list:decryptData}] {data/code/data_decrypt.java}

\section{Dart编程语言实现}
RRS 是一个Trojan和JLS的具体实现。详见地址：

https://github.com/JimmyHuang454/RRS/tree/master/lib/transport/jls

\section{节点分享格式}
示例：
\lstinputlisting[caption=分享格式,label={list:share}] {data/code/share.json}


% !TeX root = ../thuthesis-example.tex
\chapter{Specification V3}
第3个版本跟TLS完全一模一样，除了random字段是特殊生成的。

\section{基本原理}
因为TLS1.3会生成KeyShare来保证此次对话是唯一的。因为random生成时加入了KeyShare，所以攻击者是无法伪装成Client或Server。

\section{random生成算法}
生成一个随机数N（16字节），然后使用AES\_256 （随机数长度64字节）对其加密。加密时，必须加上各端 KeyShare 的内容，也就说必须要支持TLS1.3。

\section{Server 验证}
Server 根据收到Client Hello的random来判断是否有效Client，如果不是有效Client，直接转发到伪装站。如果是有效Client，则不需转发到伪装站，使用自签证书，完整按照TLS的流程处理。

\section{Client 验证}
Client 根据收到的 Server Hello的random来判断是否有效Server，如果不是有效Server，则应表现成一个正常HTTP请求。如果是有效Server，则不需要验证证书是否有效（直接信任），最后按照正常TLS流程处理即可。

% !TeX root = ../thuthesis-example.tex

\chapter{引用文献的标注}

模板支持 BibTeX 和 BibLaTeX 两种方式处理参考文献。
下文主要介绍 BibTeX 配合 \pkg{natbib} 宏包的主要使用方法。


\section{顺序编码制}

在顺序编码制下，默认的 \cs{cite} 命令同 \cs{citep} 一样，序号置于方括号中，
引文页码会放在括号外。
统一处引用的连续序号会自动用短横线连接。

\thusetup{
  cite-style = super,
}
\noindent
\begin{tabular}{l@{\quad$\Rightarrow$\quad}l}
  \verb|\cite{zhangkun1994}|               & \cite{zhangkun1994}               \\
  \verb|\citet{zhangkun1994}|              & \citet{zhangkun1994}              \\
  \verb|\citep{zhangkun1994}|              & \citep{zhangkun1994}              \\
  \verb|\cite[42]{zhangkun1994}|           & \cite[42]{zhangkun1994}           \\
  \verb|\cite{zhangkun1994,zhukezhen1973}| & \cite{zhangkun1994,zhukezhen1973} \\
\end{tabular}


也可以取消上标格式，将数字序号作为文字的一部分。
建议全文统一使用相同的格式。

\thusetup{
  cite-style = inline,
}
\noindent
\begin{tabular}{l@{\quad$\Rightarrow$\quad}l}
  \verb|\cite{zhangkun1994}|               & \cite{zhangkun1994}               \\
  \verb|\citet{zhangkun1994}|              & \citet{zhangkun1994}              \\
  \verb|\citep{zhangkun1994}|              & \citep{zhangkun1994}              \\
  \verb|\cite[42]{zhangkun1994}|           & \cite[42]{zhangkun1994}           \\
  \verb|\cite{zhangkun1994,zhukezhen1973}| & \cite{zhangkun1994,zhukezhen1973} \\
\end{tabular}



\section{著者-出版年制}

著者-出版年制下的 \cs{cite} 跟 \cs{citet} 一样。

\thusetup{
  cite-style = author-year,
}
\noindent
\begin{tabular}{@{}l@{$\Rightarrow$}l@{}}
  \verb|\cite{zhangkun1994}|                & \cite{zhangkun1994}                \\
  \verb|\citet{zhangkun1994}|               & \citet{zhangkun1994}               \\
  \verb|\citep{zhangkun1994}|               & \citep{zhangkun1994}               \\
  \verb|\cite[42]{zhangkun1994}|            & \cite[42]{zhangkun1994}            \\
  \verb|\citep{zhangkun1994,zhukezhen1973}| & \citep{zhangkun1994,zhukezhen1973} \\
\end{tabular}

\vskip 2ex
\thusetup{
  cite-style = super,
}
注意，引文参考文献的每条都要在正文中标注
\cite{zhangkun1994,zhukezhen1973,dupont1974bone,zhengkaiqing1987,%
  jiangxizhou1980,jianduju1994,merkt1995rotational,mellinger1996laser,%
  bixon1996dynamics,mahui1995,carlson1981two,taylor1983scanning,%
  taylor1981study,shimizu1983laser,atkinson1982experimental,%
  kusch1975perturbations,guangxi1993,huosini1989guwu,wangfuzhi1865songlun,%
  zhaoyaodong1998xinshidai,biaozhunhua2002tushu,chubanzhuanye2004,%
  who1970factors,peebles2001probability,baishunong1998zhiwu,%
  weinstein1974pathogenic,hanjiren1985lun,dizhi1936dizhi,%
  tushuguan1957tushuguanxue,aaas1883science,fugang2000fengsha,%
  xiaoyu2001chubanye,oclc2000about,scitor2000project%
}。



% 其他部分
\backmatter

% 参考文献
\bibliography{ref/refs}  % 参考文献使用 BibTeX 编译
% \printbibliography       % 参考文献使用 BibLaTeX 编译

% 附录
% 本科生需要将附录放到声明之后，个人简历之前
\appendix
% % !TeX root = ../thuthesis-example.tex

\begin{survey}
\label{cha:survey}

\title{Title of the Survey}
\maketitle


\tableofcontents


本科生的外文资料调研阅读报告。


\section{Figures and Tables}

\subsection{Figures}

An example figure in appendix (Figure~\ref{fig:appendix-survey-figure}).

\begin{figure}
  \centering
  \includegraphics[width=0.6\linewidth]{example-image-a.pdf}
  \caption{Example figure in appendix}
  \label{fig:appendix-survey-figure}
\end{figure}


\subsection{Tables}

An example table in appendix (Table~\ref{tab:appendix-survey-table}).

\begin{table}
  \centering
  \caption{Example table in appendix}
  \begin{tabular}{ll}
    \toprule
    File name       & Description                                         \\
    \midrule
    thuthesis.dtx   & The source file including documentaion and comments \\
    thuthesis.cls   & The template file                                   \\
    thuthesis-*.bst & BibTeX styles                                       \\
    thuthesis-*.bbx & BibLaTeX styles for bibliographies                  \\
    thuthesis-*.cbx & BibLaTeX styles for citations                       \\
    \bottomrule
  \end{tabular}
  \label{tab:appendix-survey-table}
\end{table}


\section{Equations}

An example equation in appendix (Equation~\eqref{eq:appendix-survey-equation}).
\begin{equation}
  \frac{1}{2 \uppi \symup{i}} \int_\gamma f = \sum_{k=1}^m n(\gamma; a_k) \mathscr{R}(f; a_k)
  \label{eq:appendix-survey-equation}
\end{equation}


\section{Citations}

Example citations in appendix.
\cite{abrahams99tex}
\cite{salomon1995advanced}
\cite{abrahams99tex,salomon1995advanced}


\bibliographystyle{unsrtnat}
\bibliography{ref/appendix}

\end{survey}
       % 本科生：外文资料的调研阅读报告
% % !TeX root = ../thuthesis-example.tex

\begin{translation}
\label{cha:translation}

\title{书面翻译题目}
\maketitle

\tableofcontents


本科生的外文资料书面翻译。


\section{图表示例}

\subsection{图}

附录中的图片示例（图~\ref{fig:appendix-translation-figure}）。

\begin{figure}
  \centering
  \includegraphics[width=0.6\linewidth]{example-image-a.pdf}
  \caption{附录中的图片示例}
  \label{fig:appendix-translation-figure}
\end{figure}


\subsection{表格}

附录中的表格示例（表~\ref{tab:appendix-translation-table}）。

\begin{table}
  \centering
  \caption{附录中的表格示例}
  \begin{tabular}{ll}
    \toprule
    文件名          & 描述                         \\
    \midrule
    thuthesis.dtx   & 模板的源文件，包括文档和注释 \\
    thuthesis.cls   & 模板文件                     \\
    thuthesis-*.bst & BibTeX 参考文献表样式文件    \\
    thuthesis-*.bbx & BibLaTeX 参考文献表样式文件  \\
    thuthesis-*.cbx & BibLaTeX 引用样式文件        \\
    \bottomrule
  \end{tabular}
  \label{tab:appendix-translation-table}
\end{table}


\section{数学公式}

附录中的数学公式示例（公式\eqref{eq:appendix-translation-equation}）。
\begin{equation}
  \frac{1}{2 \uppi \symup{i}} \int_\gamma f = \sum_{k=1}^m n(\gamma; a_k) \mathscr{R}(f; a_k)
  \label{eq:appendix-translation-equation}
\end{equation}


\section{文献引用}

文献引用示例\cite{abrahams99tex}。


\appendix

\section{附录}

附录的内容。


% 书面翻译的参考文献
\bibliographystyle{unsrtnat}
\bibliography{ref/appendix}

% 书面翻译对应的原文索引
\begin{translation-index}
  \nocite{salomon1995advanced}
  \bibliographystyle{unsrtnat}
  \bibliography{ref/appendix}
\end{translation-index}

\end{translation}
  % 本科生：外文资料的书面翻译
% !TeX root = ../thuthesis-example.tex

\chapter{补充内容}

附录是与论文内容密切相关、但编入正文又影响整篇论文编排的条理和逻辑性的资料，例如某些重要的数据表格、计算程序、统计表等，是论文主体的补充内容，可根据需要设置。

附录中的图、表、数学表达式、参考文献等另行编序号，与正文分开，一律用阿拉伯数字编码，
但在数码前冠以附录的序号，例如“图~\ref{fig:appendix-figure}”，
“表~\ref{tab:appendix-table}”，“式\eqref{eq:appendix-equation}”等。


\section{插图}

% 附录中的插图示例（图~\ref{fig:appendix-figure}）。

\begin{figure}
  \centering
  \includegraphics[width=0.6\linewidth]{example-image-a.pdf}
  \caption{附录中的图片示例}
  \label{fig:appendix-figure}
\end{figure}


\section{表格}

% 附录中的表格示例（表~\ref{tab:appendix-table}）。

\begin{table}
  \centering
  \caption{附录中的表格示例}
  \begin{tabular}{ll}
    \toprule
    文件名          & 描述                         \\
    \midrule
    thuthesis.dtx   & 模板的源文件，包括文档和注释 \\
    thuthesis.cls   & 模板文件                     \\
    thuthesis-*.bst & BibTeX 参考文献表样式文件    \\
    thuthesis-*.bbx & BibLaTeX 参考文献表样式文件  \\
    thuthesis-*.cbx & BibLaTeX 引用样式文件        \\
    \bottomrule
  \end{tabular}
  \label{tab:appendix-table}
\end{table}


\section{数学表达式}

% 附录中的数学表达式示例（式\eqref{eq:appendix-equation}）。
\begin{equation}
  \frac{1}{2 \uppi \symup{i}} \int_\gamma f = \sum_{k=1}^m n(\gamma; a_k) \mathscr{R}(f; a_k)
  \label{eq:appendix-equation}
\end{equation}


\section{参考文献}

附录中的参考文献示例（\cite{carlson1981two} 和 \cite{carlson1981two,taylor1983scanning,taylor1981study}）。

\printbibliography


% 致谢
% !TeX root = ../thuthesis-example.tex

\begin{acknowledgements}
	感谢父母的养育之恩，不留余力和不计代价地支持我学习与生活。萱草生堂阶，游子行天涯。
\end{acknowledgements}


% 声明
\statement
% 将签字扫描后的声明文件 scan-statement.pdf 替换原始页面
% \statement[file=scan-statement.pdf]
% 本科生编译生成的声明页默认不加页脚，插入扫描版时再补上；
% 研究生编译生成时有页眉页脚，插入扫描版时不再重复。
% 也可以手动控制是否加页眉页脚
% \statement[page-style=empty]
% \statement[file=scan-statement.pdf, page-style=plain]

% 个人简历、在学期间完成的相关学术成果
% 本科生可以附个人简历，也可以不附个人简历
% % !TeX root = ../thuthesis-example.tex

\begin{resume}

  \section*{个人简历}

  197× 年 ×× 月 ×× 日出生于四川××县。

  1992 年 9 月考入××大学化学系××化学专业，1996 年 7 月本科毕业并获得理学学士学位。

  1996 年 9 月免试进入清华大学化学系攻读××化学博士至今。


  \section*{在学期间完成的相关学术成果}

  \subsection*{学术论文}

  \begin{achievements}
    \item Yang Y, Ren T L, Zhang L T, et al. Miniature microphone with silicon-based ferroelectric thin films[J]. Integrated Ferroelectrics, 2003, 52:229-235.
    \item 杨轶, 张宁欣, 任天令, 等. 硅基铁电微声学器件中薄膜残余应力的研究[J]. 中国机械工程, 2005, 16(14):1289-1291.
    \item 杨轶, 张宁欣, 任天令, 等. 集成铁电器件中的关键工艺研究[J]. 仪器仪表学报, 2003, 24(S4):192-193.
    \item Yang Y, Ren T L, Zhu Y P, et al. PMUTs for handwriting recognition. In press[J]. (已被Integrated Ferroelectrics录用)
  \end{achievements}


  \subsection*{专利}

  \begin{achievements}
    \item 任天令, 杨轶, 朱一平, 等. 硅基铁电微声学传感器畴极化区域控制和电极连接的方法: 中国, CN1602118A[P]. 2005-03-30.
    \item Ren T L, Yang Y, Zhu Y P, et al. Piezoelectric micro acoustic sensor based on ferroelectric materials: USA, No.11/215, 102[P]. (美国发明专利申请号.)
  \end{achievements}

\end{resume}


% 指导教师/指导小组评语
% 本科生不需要
% % !TeX root = ../thuthesis-example.tex

\begin{comments}
% \begin{comments}[name = {指导小组评语}]
% \begin{comments}[name = {Comments from Thesis Supervisor}]
% \begin{comments}[name = {Comments from Thesis Supervision Committee}]

  论文提出了……

\end{comments}


% 答辩委员会决议书
% 本科生不需要
% % !TeX root = ../thuthesis-example.tex

\begin{resolution}

  论文提出了……

  论文取得的主要创新性成果包括：

  1. ……

  2. ……

  3. ……

  论文工作表明作者在×××××具有×××××知识，具有××××能力，论文××××，答辩××××。

  答辩委员会表决，（×票/一致）同意通过论文答辩，并建议授予×××（姓名）×××（门类）学博士/硕士学位。

\end{resolution}


% 本科生的综合论文训练记录表（扫描版）
% \record{file=scan-record.pdf}

\end{document}
